% Section 3.3: Allocation (deployment) view. Owner: Slice 5 (issue #21).
\subsection{Allocation View}
\label{sec:allocation-view}

\subsubsection{View description}
We use the deployment style to allocate the runtime elements of
Section~\ref{sec:cc-view} to execution hosts, storage and communication links.
Deployment structures relate software to its hardware environment (Bass, Clements,
and Kazman 15). This packet shows one store repeated across the chain and two
\DC hosts. It identifies which failures affect local service and which durable
state supports recovery; runtime exchanges remain in the C\&C view.

\subsubsection{Primary presentation}
Figure~\ref{fig:deployment} shows the baseline deployment. Each component is nested
inside its execution host, and persistent state is nested inside that host's
storage. The WAN is inside our system boundary: it connects the store and \DC
nodes, even though another organization may operate the physical network.
The POS system and \AnalyticsClient are outside the boundary.

\begin{figure}[htbp]
\centering
\includegraphics[width=\linewidth,height=0.8\textheight,keepaspectratio]{deployment}
\caption{Allocation view --- deployment of store and distribution-center components, persistent state and network links.}
\label{fig:deployment}
\end{figure}

\subsubsection{Element catalog}
\paragraph{Nodes and links.}
Table~\ref{tab:allocation-nodes} records placements and capacity obligations.
Rates are design requirements for provisioning and load tests, not measurements
of selected hardware. We do not prescribe processor, disk or operating-system
models without those tests.

\begin{table}[htbp]
\centering
\small
\begin{tabularx}{\linewidth}{|>{\raggedright\arraybackslash}p{0.19\linewidth}|>{\raggedright\arraybackslash}X|}
\hline
\textbf{Node/link} & \textbf{Allocation and properties}\\\hline
Store node ($\times N$) & $N=3$ initially, up to 10. Hosts \POSAdapter and \StoreAgent; local storage holds \StoreLocalInventory, \Outbox and recovery history. Sustain 100 events/s locally (\scnref{PERF-1}); size buffering in days: 8{,}640{,}000 events per day, plus orders and headroom. Support 2{,}000 events/s of catch-up while live sales continue (\scnref{PERF-4}).\\\hline
Primary \DC node & Hosts \CFE, \DCAgent, \CentralMasterInventory, \ShippingOrderFulfillment, \StoreTrackingClient and \ReportingClient. Local storage holds the CFE event store, master inventory database, source/subscription positions and checkpoints. Absorb 1{,}000 live store events/s (\scnref{PERF-2}), plus \DC events, orders, replay and fan-out to both masters. Provision additional ingress/storage capacity for each concurrently recovering store's \scnref{PERF-4} burst.\\\hline
Second \DC host & Hosts \StandbyMaster with its own database, applied positions and duplicate-detection state. A second \CFE subscription applies every event idempotently. Provides the \InventoryQuery after failover; capacity must cover live delivery, queries and copying state for recovery (\scnref{AV-3}).\\\hline
\DC inter-host link & Carries \IEP delivery/acknowledgement, master state transfer and failover coordination. Network \IEP and \InventoryQuery connections use TLS 1.3 with mutual certificates, including inside the \DC; protect recovery transfer likewise.\\\hline
WAN link & Internal store-to-\DC communication path for the \IEP. TLS 1.3 with mutual certificates protects both directions (\scnref{SEC-2}). Provision for live traffic and recovery; byte-rate sizing depends on message size. Disconnection delays central updates without requiring local transactions to wait.\\\hline
External POS node & One external system per store; sales enter through the \POSAdapter as a one-way stream. Its implementation and storage remain outside our deployment.\\\hline
External analytics node & Hosts the future \AnalyticsClient outside our boundary. Uses the read-only \InventoryQuery on the active master over mutual TLS, with authorization and rate limits (\scnref{INT-1}); the same controls apply on failover.\\\hline
\end{tabularx}
\caption{Allocation-view node catalog: placements, storage and capacity properties.}
\label{tab:allocation-nodes}
\end{table}

\paragraph{Persistence and failure assumptions.}
Durable storage survives process restarts; local transactions atomically commit
inventory changes and outgoing records. The store's retained local log and a
consistent snapshot are authoritative for rebuilding \StoreLocalInventory.
Pruning acknowledged \Outbox entries does not remove that recovery history.
At the \DC, the CFE event store is authoritative for replaying accepted messages
into both master projections. Each master database commits its applied cursor with
its updates. Source/subscription positions and duplicate-detection state also
survive restarts; memory or socket receipt alone is insufficient.

A WAN outage leaves store reads and updates available while the store process,
disk and capacity survive (\scnref{AV-1}). Replay is normal for any outage length;
days of buffering are a provisioning choice, not a bound on the outage. Catch-up
time grows with the backlog; missing history uses the \StoreSnapshot recovery path.
The separate local-access requirement is at least 99\% availability and no
individual local-inventory outage longer than 30 minutes. A supervised restart
and local rebuild must restore usable access within that limit, including
detection time; \scnref{TEST-2} measures the interval. If a host cannot restart,
replacement and durable-state recovery are required; supervision alone does not
repair hardware.

A \CentralMasterInventory process failure leaves the \CFE accepting messages.
Query clients fail over within seconds to the \StandbyMaster (\scnref{AV-3}). Only
the master holding the lease serves queries: it stops if renewal fails, and the
standby takes over only after lease expiry. Lease coordination must preserve that
exclusivity during a partition. The failed copy restores a consistent copy of the
survivor's state and per-source applied positions, then replays and rejoins as standby.
Failure of the whole Primary \DC node also stops its \CFE and client services:
stores buffer locally and central acceptance and reports pause; the Second \DC host
can serve queries over its last applied state, but new updates wait for \CFE recovery.
The standby does not replicate the event bus or reporting services. Disk recovery
requires surviving state and subsequent retained history; permanent loss of every
durable copy cannot meet the no-loss guarantee. Validate lease expiry, state transfer
and recovery along with failover timing.

\paragraph{Interfaces, behavior and constraints.}
Event delivery obeys the \IEP (Section~\ref{sec:iep}); read-only queries use the
\InventoryQuery on the active master. Local persistence
uses the owning service's transactional storage interface. Each internal C\&C
component/data store has exactly one baseline placement, with one instance per
store where applicable. External systems are not allocated to our hosts.
Recovery compares per-source watermarks before replay, as shown in
Figure~\ref{fig:cc-seq-reconcile}; this packet identifies the surviving storage
used by that behavior rather than defining another recovery protocol.

\subsubsection{Context diagram}
The enclosing boundary and external nodes in Figure~\ref{fig:deployment} also
serve as this packet's context diagram. Its external interactions match
Figure~\ref{fig:cc-context}: POS sales enter the store, and the \AnalyticsClient
queries the active master through the \InventoryQuery. Manager roles use the store and \DC applications
already shown in the C\&C view; they do not introduce additional deployment hosts.
In particular, the store-to-\DC WAN is internal to the system, not a third
external system.

\subsubsection{Variability mechanisms}
\begin{itemize}[nosep]
\item \emph{Store count.} Add a store-node instance, its local storage,
\StoreAgent configuration, source identity and certificates. The component
placements repeat; capacity must cover the added traffic (\scnref{SCAL-1}).
\item \emph{\DC host count.} The primary services and hot standby occupy two
hosts. A later deployment may separate other \CFE and client services across
additional hosts, preserving their logical interfaces and storage ownership. Such a change
requires an updated mapping, secured inter-host links and capacity/failure tests;
splitting services alone does not establish redundant failover.
\end{itemize}

\subsubsection{Architecture background}
We keep store inventory and outgoing messages on store-local storage so WAN
failure does not interrupt a local transaction. The second \DC host separates the
master copies' failure boundaries, at the cost of another database and secured
replication traffic. Other primary services still share a host failure boundary
and compete for CPU, disk and network capacity. The workload
in Section~\ref{sec:analysis-results} must be tested with replay and reporting,
not just steady ingestion. We distinguish process recovery from host/storage
recovery so the deployment does not imply availability that it cannot supply.

\subsubsection{Related view packets}
Section~\ref{sec:cc-view} defines the runtime components and recovery behavior;
Section~\ref{sec:module-view} defines the implementation units that realize them.
Table~\ref{tab:component-node-mapping} in Section~\ref{sec:relations} collects
the component-to-node placements. The same component
identities are used here; local recovery logs, checkpoints and subscription state
are supporting persistent state rather than additional application components.
